% config/commands.tex — Author-facing macros: notation, difficulty, tags
% -----------------------------------------------------------------------
% Requires: xcolor, amsmath, amssymb, tikz (shapes.geometric) — all
% loaded in main.tex before this file.

% =====================  Difficulty indicators  =========================
% \difficulty{n} renders n filled stars + (5-n) empty stars, gold.
% Convention: * = routine, ** = standard olympiad, *** = hard,
%             **** = very hard (national team), ***** = IMO P3/P6 level.
\usetikzlibrary{shapes.geometric}
% One inline star; #1 = fill color, #2 = draw color.
\newcommand{\staricon}[2]{%
  \begin{tikzpicture}[baseline=(s.center)]
    \node[star, star points=5, star point ratio=2.3, inner sep=0.85pt,
          fill=#1, draw=#2, line width=0.45pt] (s) {};
  \end{tikzpicture}%
}
\newcommand{\filledstars}[1]{% #1 filled stars remain
  \ifnum#1>0 \staricon{gold}{gold}\expandafter\filledstars\expandafter{\the\numexpr#1-1}\fi
}
\newcommand{\emptystars}[1]{% #1 outline stars remain
  \ifnum#1>0 \staricon{white}{rulegray}\expandafter\emptystars\expandafter{\the\numexpr#1-1}\fi
}
\newcommand{\difficulty}[1]{\filledstars{#1}\emptystars{\numexpr5-#1\relax}}

% Variant for use on colored title bars (e.g. gold challenge bars):
% white filled stars, pale-gold outline stars.
\newcommand{\filledstarsW}[1]{%
  \ifnum#1>0 \staricon{white}{white}\expandafter\filledstarsW\expandafter{\the\numexpr#1-1}\fi
}
\newcommand{\emptystarsW}[1]{%
  \ifnum#1>0 \staricon{white}{gold!45}\expandafter\emptystarsW\expandafter{\the\numexpr#1-1}\fi
}
\newcommand{\difficultywhite}[1]{\filledstarsW{#1}\emptystarsW{\numexpr5-#1\relax}}

% Text difficulty badge for prose: \difficultytext{2} -> "Level 2"
\newcommand{\difficultytext}[1]{\textbf{\color{gold}Level #1}}

% =====================  Problem metadata  ==============================
% \source{...} — right-aligned provenance tag, e.g. \source{IMO 2005, P1}
\newcommand{\source}[1]{\hfill{\small\sffamily\itshape\color{slate}[#1]}\par}

% \points{n} — optional point value tag for assessments.
\newcommand{\points}[1]{\hfill{\small\sffamily\color{slate}(#1~pt)}\par}

% =====================  Number sets  ====================================
\newcommand{\N}{\mathbb{N}}
\newcommand{\Z}{\mathbb{Z}}
\newcommand{\Q}{\mathbb{Q}}
\newcommand{\R}{\mathbb{R}}
\newcommand{\C}{\mathbb{C}}
\newcommand{\Positive}{\R_{>0}}   % strictly positive reals

% =====================  General math shortcuts  ========================
\DeclareMathOperator{\sgn}{sgn}
\newcommand{\abs}[1]{\lvert #1 \rvert}
\newcommand{\norm}[1]{\lVert #1 \rVert}
\newcommand{\floor}[1]{\left\lfloor #1 \right\rfloor}
\newcommand{\ceil}[1]{\left\lceil #1 \right\rceil}
% Display-style fractions and sums for inline use:
\newcommand{\ds}{\displaystyle}
% Standard big operators with canonical limits:
\newcommand{\sumi}{\sum_{i=1}^{n}}
\newcommand{\sumk}{\sum_{k=1}^{n}}
% QED-flavored end-of-proof symbol is provided by amsthm (\qedhere).

% =====================  Curriculum metadata  ===========================
% Overridden per project (see scripts/init-curriculum.sh).
\newcommand{\CurriculumSeries}{Mathematics Olympiad Series}
\newcommand{\CurriculumTitle}{Curriculum Title}
\newcommand{\CurriculumSubtitle}{Subtitle: Volume and Level}
\newcommand{\CurriculumAuthor}{Author Name}
\newcommand{\CurriculumInstitution}{Institution}
\newcommand{\CurriculumDate}{\today}
\newcommand{\CurriculumEdition}{Student Edition}
\newcommand{\CurriculumVersion}{v0.1.0}
